{\bf Setup.}
\begin{itemize}\itemsep=0pt
	\item We adopt the following notation for the vertex operator corresponding to the state $a$
\[
Y (a, z) = \sum _{i \in \mathbb{Z}} a_{(i)}z^{-i-1}.
\]
\item The Zhu algebra associated to the vertex operator algebra $V$ with the Virasoro vector $\omega$ will be denoted with $A_{\omega}(V)$.
\item Let $\F_{\pm 1}$ be the lattice vertex {super}algebras associated to the lattice ${\mathbb Z}\sqrt{\pm 1}$ defined in Section \ref{F_{-1}}.
\item Let $M_{\alpha}(1)$ denote the rank one Heisenberg vertex algebra generated by a Heisenberg vector~$\alpha$.
\item Let $\mathcal{W}^k(\mathfrak{sl}_n, f_{\rm sub})$ be the (universal) subregular $\mathcal{W}$-algebra corresponding to the subregular nilpotent element $f_{\rm sub}$ and $\mathfrak g = \mathfrak{sl}_n$, and $\mathcal{W}_k(\mathfrak{sl}_n, f_{\rm sub})$ its simple quotient.
\item $L_s(\mathfrak{sl}_2)$ is affine Lie algebra of $\mathfrak{sl}_2$ and $N_s(\mathfrak{sl}_2)$ its parafermion (coset) subalgebra.
\item The universal $N=2$ superconformal vertex algebra of central charge $c$ is denoted by $V^{N=2}_c$ and $L^{N=2}_{c}$ is its simple quotient.
\end{itemize}

\section{Preliminaries}	

\subsection[Lattice vertex superalgebras F\_1]{ Lattice vertex superalgebras $\boldsymbol{\mathcal F_{\pm 1}}$} \label{F_{-1}}

Consider a rank one lattice $L= {\Z}\varphi^{\pm} $, $\big\langle \varphi^{\pm}, \varphi^{\pm} \big\rangle = \pm 1$. Let $\F_{1}$ (resp.\ $\F_{-1}$) be the associated vertex algebra. These vertex superalgebras are used for the construction of the inverse of Kazama--Suzuki functor in the context of duality between affine $\widehat{sl}_2$ and $N=2$ superconformal algebra (cf.\ \cite{A-IMRN,A-2001,FST}).

As a vector space $\F_{\pm 1} = {\C}[L] \otimes M_{\varphi^{\pm}} (1)$, where ${\C}[L]$ is a group algebra of $L$, and $M_{\varphi^{\pm}} (1)$ the Heisenberg vertex algebra generated by the Heisenberg field $\varphi^{\pm}(z) = \sum_{n \in {\Z} } \varphi^{\pm}(n) z^{-n-1}$ such that
\[
 \bigl[\varphi^{\pm}(n), \varphi^{\pm}(m)\bigr] = \pm n \delta_{n+m, 0}.
 \]

The vertex algebras $\F_{\pm 1}$ are weakly generated by \smash{$e^{ \varphi^{\pm}}$} and \smash{$e^{-\varphi^{\pm}}$}.
Moreover, $\F_{\pm 1}$ is a simple vertex superalgebra and a completely reducible $M_{\varphi^{\pm}} (1)$-module isomorphic to
\[
\F_{\pm 1} \cong \bigoplus_{m \in \mathbb{Z}} \F_{\pm 1}^m ,
\]
where $ \F_{\pm 1}^m$ is an irreducible $M_{\varphi^{\pm}} (1)$-module generated by \smash{$e^{m\varphi^{\pm}}$}.


\subsection[N=2 superconformal algebra]{$\boldsymbol{ N=2}$ superconformal algebra}

The universal $N=2$ superconformal vertex algebra $V^{N=2}_{c}$ is generated by even fields $T(z)$, $H(z)$ and odd fields $E(z)$, $F(z)$
satisfying the following OPEs
\begin{gather*}
 T(z)T(w) \sim \frac{\frac{c}{2}}{(z-w)^{4}}+\frac{2T(w)}{(z-w)^{2}} + \frac{\partial T(w)}{z-w}, \qquad H(x)H(y) \sim \frac{ \frac{c}{3}}{(z-w)^{2}}, \qquad \\
 T(z)H(w) \sim \frac{H(w)}{ (z-w)^{2}}+ \frac{\partial H(w)}{z-w}, \\
 T(z)E(w) \sim \frac{\frac{3}{2} E(w)}{(z-w)^{2}} + \frac{\partial E(w)}{z-w}, \qquad T(z)F(w) \sim \frac{\frac{3}{2} F(w)}{(z-w)^{2}} + \frac{\partial F(w)}{z-w}, \\
	E(z)F(w) \sim \frac{\frac{2c}{3}}{(z-w)^{3}}+\frac{2H(w)}{(z-w)^{2}} + \frac{2T(w) + \partial H(w)}{z-w}, \\
	F(z)E(w) \sim \frac{\frac{2c}{3}}{(z-w)^{3}}-\frac{2H(w)}{(z-w)^{2}} + \frac{2T(w) - \partial H(w)}{z-w}, \\
	H(z)E(w) \sim \frac{E(w)}{z-w}, \qquad H(z)F(w) \sim -\frac{F(w)}{z-w}, \qquad E(z)E(w) \sim 0, \qquad F(z)F(w) \sim 0. 	
\end{gather*}

Set
\begin{gather*}
 T(z) = \sum _{i \in {\Z}} T(i) z^{-i-2}, \qquad H(z) = \sum _{i \in {\Z}} H(i) z^{-i-1},\\
 E(z) = \sum _{i \in {\Z}} E_{(i)} z^{-i-1} = \sum _{i \in {\Z}} E\left(i+\frac{1}{2}\right)z^{-i-2},\\
 F(z) = \sum _{i \in {\Z}} F_{(i)} z^{-i-1} = \sum _{i \in {\Z}} F\left(i+\frac{1}{2}\right)z^{-i-2}.
 \end{gather*}
 Then the components of these fields satisfy the commutation relation for the $N=2$ superconformal algebra with basis $\{T(n), H(n), E(r), F(r)\}$, $n \in \mathbb{Z}$, $r \in \frac{1}{2}+\mathbb{Z}$,
\begin{gather*}
	[T(m), T(n)] = (m-n)T(m+n) + \frac{c}{12}\bigl(m^3 - m\bigr)\delta_{m+n,0}, \qquad [H(m), H(n)] = \frac{c}{3}m\delta_{m+n,0}, \\
	[T(m), H(n)] = -nH(n+m), \qquad [ T(m) , E(r) ] = \left(\frac{1}{2} m-r \right)E(m+r), \\
 [ T(m) , F(r) ] = \left(\frac{1}{2} m-r \right)F(m+r), \\
	[ H(m) ,E(r)] = E(m+r), \qquad [ H(m) , F(r)] = -F(m+r), \\
	\{E(r), F(s) \} = 2 T({r + s}) + ( r - s )H(r + s) + \frac{c}{3}\left( r^2 - \frac{1}{4} \right)\delta_{r+s,0}, \\
	\{E(r), E(s) \} = \{F(r), F(s) \} = 0.
\end{gather*}

 Let $L^{N=2}_{c}$ be the simple quotient of $V^{N=2}_{c}$.

\subsubsection{Spectral flow automorphisms}
The $N=2$ superconformal algebra $V^{N=2}_{c}$ admits a family of spectral flow automorphisms $\sigma^\ell$, $\ell \in \mathbb{Z}$, given by
\begin{align*}
\begin{aligned}
	& \sigma^\ell(T(n))= T(n) -\ell H(n) + \frac{1}{6}\ell ^2\delta_{n,0}c\mathbbm{1}, \qquad \sigma^l(H(n)) = H(n) + \frac{1}{3}\ell \delta_{n,0}c\mathbbm{1},  \\
	& \sigma^\ell(E(r))= E({r-\ell}), \qquad \sigma^\ell(F(r)) = F({r+\ell}), \qquad \sigma^\ell(\mathbbm{1}) = \mathbbm{1}.
\end{aligned}
\end{align*}



Define
\[
 \Delta (-H,z) = z^{-H(0)}\exp \left ( \sum _{k=1} ^{\infty} (-1)^{k+1}\frac{-H(k) }{kz^k}\right ).
 \]

For any $V^{N=2}_{c}$-module $(M, Y_M(\cdot, z))$ and $n \in \mathbb{Z}$,
$\sigma^ n (M)$ is again a $L^{N=2}_{c}$-module with vertex operator structure given by
$ Y_{\sigma ^n (M) } ( \cdot , z)) := Y_{ M } (\Delta ( -n H, z) \cdot , z))$.

\subsubsection{Twisted highest weight conditions}
Let us define a new Virasoro vector $\widetilde T = T + \partial H$. Then $E(z)$ is primary field of conformal weight~$1/2$ and $F(z)$ of conformal weight $5/2$ with respect to $\widetilde T$.
 With respect to the new grading operator $\widetilde T(0)$ we shall also need to introduce twisted highest weight conditions.
 Let~${(h,q) \in {\C}^2}$. Let $M_c[h,q]$ be the Verma module of the $N=2$ superconformal algebra of central charge $c$ generated by the (twisted) highest weight vector $v_{h,q}$ such that for $n \in {\Z}_{\ge 0}$ (cf.\ \cite{FST}):
\begin{gather*}
 H(n) v_{h,q} =\delta_{n, 0} h v_{h,q}, \qquad T(n) v_{h,q} = \delta_{n,0} q v_{h,q}, \\ E\left(n-\frac{1}{2}\right) v_{h,q} = F\left(n+\frac{3}{2}\right) v_{h,q} =0.
 \end{gather*}
 Let $L_c[h,q]$ be the irreducible quotient of $M_c[h,q]$. Then $L_c[h,q]$ is an irreducible $V_c^{N=2}$-module.

\subsubsection{Parafermionic subalgebras}
\label{paraferm-1}
Let $ \mathcal N_c = \operatorname{Com} \bigl(M_H(1) , L^{N=2}_{c}\bigr) $ be the parafermionic subalgebra of $L^{N=2}_{c}$. Note that using the Kazama--Suzuki duality between $L^{N=2} _c$ and $L_k(\mathfrak{sl}_2)$ (cf.\ \cite{A-IMRN, FST}), we have that
$\mathcal N_c = \mathcal N_s (\mathfrak{sl}_2)$, where $c =\frac{3s}{s+2}$ and
$\mathcal N_s (\mathfrak{sl}_2)$ is the parafermion vertex subalgebra of $L_s(\mathfrak{sl}_2)$ (cf.\ \cite{DLY}). If $c \notin \{ 0, 1, \frac{3}{2}, -6, -9\}$, then the vertex algebra $ \mathcal N_c $ is (weakly) generated by the fields $ T^{\perp}$ and~$W^{N=2}_{c}$ (cf.\ \cite{BEHHH,DLY}), where
\begin{gather*}
	 T^{\perp} = T - \frac{3}{2c}{:}HH{:}, \qquad
	W^{N=2}_{c} = \nu \left({:} E F{:} - \partial T -\frac{6}{c}{:} T H{:} - \frac{c-9}{3c}\partial^2H + \frac{6}{c^2}{:}H^3{:} \right),
\end{gather*}
and $\nu \in \mathbb{C}$ is a normalization factor. In particular, $\mathcal N_{c=-15}$ is weakly generated by $T^{\perp}$ and $W^{N=2}_{c=-15}$.

 For other central charges, we have the following:
\begin{itemize}\itemsep=0pt
\item If $c\in \{ 0, 1\}$, then $\mathcal N_c = {\C} {\bf 1}$ (cf.\ \cite{DLY}).
\item $N_{c=3/2}$ is the simple Virasoro vertex algebra of central charge $c=\frac{1}{2}$.
 \item \smash{$\mathcal N_{c=-6}=\mathcal N_{-4/3}(\mathfrak{sl}_2)$} is the singlet vertex algebra $\mathcal M(3)$ generated by the Virasoro field and primary field of conformal weight $5$ (cf.\ \cite{A-2005}).
 \item \smash{$\mathcal N_{c=-9} =\mathcal N_{-3/2} (\mathfrak{sl}_2)$} is a direct sum of \smash{$\mathcal W_{-5/2} (\mathfrak{sl}_3, f_{\rm pr})$}-modules (cf.\ \cite{AMW}). As a vertex algebra it is generated by \smash{$\mathcal W_{-5/2} (\mathfrak{sl}_3, f_{\rm pr})$} and its simple module of conformal weight $4$.
\end{itemize}

 Let $A_{\omega}(V)$ be the Zhu algebra associated to the VOA $V$ with the Virasoro vector $\omega$, and let~$[v]$ be the image of $v \in V $ under the mapping $V \mapsto A_{\omega}(V)$.

 One can show that the Zhu's algebra $A\bigl(V^{N=2}_{c}\bigr)$ (cf.\ \cite{A-IMRN}) is isomorphic to ${\mathbb C}[x,y]$; i.e., there is isomorphism $f \colon A\bigl(V^{N=2}_{c}\bigr) \rightarrow {\mathbb C}[x_1,y_1]$ such that
$[H({-1}){\mathbbm 1}] \mapsto x_1, \ [T({-2}{)\mathbbm 1}] \mapsto y_1$.

\subsection{Kazama--Suzuki duality}
In this subsection, we will define a duality of vertex algebras which is motivated by the duality between $N=2$ superconformal vertex algebra and affine vertex algebra $L_k(\mathfrak{sl}_2)$.

Recall first that if $S$ is a vertex subalgebra of $V$, we have the commutant subalgebra of $V$ (cf.\ \cite{LL})
\[
\operatorname{Com} (S, V):= \{ v \in V \mid a_n v = 0,\, \forall a \in S,\, \forall n \in {\Z}_{\ge 0}\}.
\]

Assume that $U$, $V$ are vertex superalgebras. We say that $V$ is the Kazama--Suzuki dual of~$U$ if there exist injective homomorphisms of vertex superalgebras
$
 \varphi_1 \colon V \rightarrow U \otimes \F_1$, $ \varphi_2 \colon U \rightarrow V \otimes \F_{-1}$,
so that
$
V \cong \operatorname{Com} \left( M_{H_1} (1), U \otimes \F_1\right)$, $ U \cong \operatorname{Com} \left( M_{H_2} (1), V \otimes \F_{-1}\right)$,
where $M_{H_1} (1) $ (resp.\ $ M_{H_2} (1)$) is a rank one Heisenberg vertex subalgebra of
$U\otimes \F_1$ (resp.\ $V\otimes \F_{-1}$), and $\F_{\pm 1}$ are lattice vertex superalgebras defined in Section~\ref{F_{-1}}.

\subsection[Subregular W-algebras W\^{}k(sl\_4, f\_sub]{Subregular $\boldsymbol{\mathcal{W}}$-algebras $\boldsymbol{\mathcal{W}^k(\mathfrak{sl}_4, f_{\rm sub})}$}	\label{subreg_ope}

Let $\mathfrak{g}$ be a simple finite-dimensional Lie algebra and $k \in \mathbb{C}$. Let $\mathcal{W}^k(\mathfrak{g}, f_{\rm sub})$ be the (universal) subregular $\mathcal{W}$-algebra corresponding to the subregular nilpotent element $f_{\rm sub}$ (\cite{KW}).

In the case where $\mathfrak g = \mathfrak{sl}_n$, it was proved in \cite{Gen} that $\mathcal{W}^k(\mathfrak{g}, f_{\rm sub})$ is isomorphic to the Feigin--Semikhatov algebra \smash{$W_n^{(2)}$} (cf.\ \cite{FS}), using a certain free field realization of $\mathcal{W}^k(\mathfrak{g}, f_{\rm sub})$.

In this paper, we will only consider the case $\mathfrak g = \mathfrak{sl}_4$. The vertex algebra $\mathcal{W}^{k}(\mathfrak{sl}_4, f_{\rm sub})$ has central charge
\[
c_k = -\frac{(3k+8)(8k+17)}{(k+4) },
\]
 and it is freely generated by fields $J(z)$, $L(z)$, $G^+ (z)$, $G^- (z)$, $W(z)$. The explicit OPEs are known (see \cite{CL18,FS}) and are written in Appendix~\ref{appendix-b}. In this paper, we shall assume that the gradation on~$\mathcal{W}^{k}(\mathfrak{sl}_4, f_{\rm sub})$ is defined by the shifted Virasoro field $\widetilde L = L + \partial J$. Then $G^+$ (resp.~$G^-$) is a~primary field of conformal weight $1$ (resp.~$3$). We set
\begin{gather*}
J(z) = \sum _{n \in {\Z}} J(n) z^{-n-1}, \qquad L(z) = \sum _{n \in {\Z}} L(n) z^{-n-2}, \qquad G^+ (z) = \sum _{n \in {\Z}} G^+(n) z^{-n-1}, \\ G^- (z) = \sum _{n \in {\Z}} G^-(n) z^{-n-3}, \qquad W (z) = \sum _{n \in {\Z}} W(n) z^{-n-3}.
\end{gather*}	
The Zhu algebra $A\bigl(\mathcal{W}^{k}(\mathfrak{sl}_4, f_{\rm sub})\bigr)$ is generated by
\[
 E=[G^+],\qquad F=[G^- ], \qquad X=[J ], \qquad Y =[L], \qquad Z=[W].
 \]

The subregular $\mathcal{W}$-algebra $\mathcal{W}^k(\mathfrak{sl}_4, f_{\rm sub})$ admits a family of spectral flow automorphisms $\psi^m$, $m \in \mathbb{Z}$, given by
\begin{gather*}
 \psi^m (J(n)) = J(n)- { \frac{m (3 k +8)}{4} }
	 \delta _{n,0}\mathbbm{1} , \qquad \psi^m ( L(n)) = L(n) - mJ(n) + \frac{m^2 { (3k+8)}}{{8} }\delta _{n,0}\mathbbm{1}, \\
 \psi^m (G^+(n)) = G^+(n-m), \qquad \psi^m (G^-(n)) = G^-(n+m).
 \end{gather*}

Define
 \[
 \Delta (-J,z) = z^{-J(0)}\exp \left ( \sum _{k=1} ^{\infty} (-1)^{k+1}\frac{-J(k)}{kz^k}\right ).
 \]

For any $\mathcal{W}^k(\mathfrak{sl}_4, f_{\rm sub})$-module $(M, Y_M(\cdot, z))$ and $n \in \mathbb{Z}$,
$\psi^m (M)$ is again a $\mathcal{W}^k(\mathfrak{sl}_4, f_{\rm sub})$-module with vertex operator structure given by
$ Y_{\psi^m (M) } ( \cdot , z)) := Y_{ M } (\Delta ( -m J, z) \cdot , z))$.

\subsection[Singular vectors in W\^{}k(sl\_n, f\_sub)]{Singular vectors in $\boldsymbol{\mathcal{W}^k(\mathfrak{sl}_n, f_{\rm sub})}$}
	
	Formulas for singular vectors in $\mathcal{W}^k(\mathfrak{sl}_n, f_{\rm sub})$ are not known in general; however, the following criterion from \cite{Feh} tells us when $(G^+)^s$ is singular.
	
	\begin{Proposition}[\cite{Feh}]\label{kriterij}
		The vector $(G^+)^s$, $s>0$, is singular in $\mathcal{W}^k(\mathfrak{sl}_n, f_{\rm sub})$ if and only if~${
		 	i(k+n-1)=s}
$
 for some $i \in \{1,\dots ,n-1 \}$.
	\end{Proposition}

Using the above criterion, it is easy to observe the following.
	\begin{Lemma} \label{sing_G^+}\qquad
	\begin{itemize}\itemsep=0pt
		\item[$(i)$] For $k=-n+3$, $(G^+)^2$ is singular in $\mathcal{W}^k(\mathfrak{sl}_n, f_{\rm sub})$ but $G^+$ is not.
		\item[$(ii)$] For $k=-n+4$, $(G^+)^3$ is singular in $\mathcal{W}^k(\mathfrak{sl}_n, f_{\rm sub})$ but $(G^+)^2$ is not.	
	\end{itemize}
	\end{Lemma}
	
	\begin{proof}
		(i) Assume that $(G^+)^2$ is singular in $\mathcal{W}^k(\mathfrak{sl}_n), f_{\rm sub})$ but $G^+$ is not. Then from Proposition \ref{kriterij} it follows that there exists $i \in \{1,\dots ,n-1 \}$ such that $ i(k+n-1)=2 $ and there does not exist $i \in \{1,\dots ,n-1 \}$ such that $ i(k+n-1)=1$. Let now $k=-n+3$. Then $i=1$ is solution of the first equation, and the second equation does not have solutions. This proves (i).
		Analogous reasoning yields (ii).
	\end{proof}
